\section{Reversing the first-pole transition}
\label{us:sec:inverse}

The forward transition answers how close an alternating harmonic sum
is to its first term. The inverse question is equally natural: how
large must the real outer exponent be to attain a specified normalized
fraction? The monotonicity needed to pose this question is exact.

\begin{lemma}[Strict order in the outer exponent]
\label{us:lem:monotone}
For fixed $h\ge1$ and $a>0$, the function $p\mapsto R_{p,h}(a)$
is continuous and strictly increasing on $(0,\infty)$, with endpoint
limits $(\log2)^h/2$ and $1$.
\end{lemma}
\begin{proof}
For $d>0$, independent gamma variables with common rate $h+a$
can be coupled by $Y_{p+d}=Y_p+Y_d$. Since
$y\mapsto k_h(e^{-y})$ is strictly increasing and $Y_d>0$
almost surely, \eqref{us:eq:gamma} proves strict increase.
As $p\downarrow0$, $Y_p\to0$ in probability, for instance by its
mean. As $p\to\infty$, $Y_p\to\infty$ in probability by its
mean and variance. The bounded continuous kernel gives the two
limits. The same gamma representation, or dominated convergence
in \eqref{us:eq:integral} on compact $p$-intervals, gives continuity.
\end{proof}

For $0<\theta<1$ put
\begin{equation}\label{us:eq:inverse-parameters}
 \lambda=-\log\theta,
 \qquad L_0=\log h-\log(2\lambda).
\end{equation}
For all sufficiently large $h$, Lemma~\ref{us:lem:monotone}
therefore gives a unique $p_\theta(h,a)>0$ satisfying
$R_{p_\theta,h}(a)=\theta$.

\begin{theorem}[All-orders inverse transition]
\label{us:thm:inverse}
Fix compact intervals of $a$ in $(0,\infty)$ and of $\theta$
in $(0,1)$. For every fixed integer $J\ge0$,
\begin{equation}\label{us:eq:inverse-all}
 p_\theta(h,a)=hL_0+
   \sum_{j=0}^J\frac{Q_j(L_0;\lambda,a)}{h^j}
   +O_J\!\left(\frac{L_0^{J+2}}{h^{J+1}}\right),
 \qquad h\to\infty,
\end{equation}
uniformly on those intervals. Here $Q_j\in\mathbb Q[L_0,\lambda,a]$
has degree at most $j+1$ in $L_0$. The first two polynomials are
\begin{align}
 Q_0(L;\lambda,a)
  & =\left(a+\frac12-\frac\lambda2\right)L
       +2-\frac{5\lambda}{6},\label{us:eq:Q0}\\
 Q_1(L;\lambda,a)
  & =\left(\frac{\lambda^2}{8}-\frac\lambda4\right)L^2
    +\left(\frac{\lambda^2}{3}-\frac\lambda3-\frac1{12}\right)L
       \notag\\
  &\hspace{1.5em}-\frac{5a\lambda}{6}+2a
       +\frac{3\lambda^2}{8}-\frac\lambda{12}-1.
       \label{us:eq:Q1}
\end{align}
\end{theorem}

\subsection{The forward coefficients used in the inversion}

We state the input with its provenance, since it is already proved in
the baseline. Write $L=p/h$ and $v=he^{-L}/2$. In the compact
transition window $L=\log h+O(1)$, the all-orders compiler of
\cite[\src{chapters/04-depth.tex}, Theorem
\texttt{harmonic:thm:allorders}]{cont:ProveIt} gives
\begin{equation}\label{us:eq:forward}
 R_{p,h}(a)=e^{-v}
    \sum_{j=0}^N\frac{\Theta_j(L,v;a)}{h^j}
       +O_N\!\left(\frac{L^{N+1}}{h^{N+1}}\right).
\end{equation}
Each $\Theta_j$ has degree at most $j$ in $L$, and
\begin{equation}\label{us:eq:Theta1}
 \Theta_0=1,
 \qquad
 \Theta_1=\left(\frac{v^2}{2}-(a+\tfrac12)v\right)L
              +\frac56v^2-2v.
\end{equation}
A finite way to regenerate these polynomials is as follows. Set
\[
 \mathcal D P(v)=vP(v)-vP'(v),\qquad
 d_j(t;a)=\frac{(-1)^j}{j+1}
                \{a^{j+1}-(a-t)^{j+1}\}.
\]
Define $\mathcal U_j$ by the formal kernel expansion
\[
 k_h(2v/h)=e^{-v}\sum_{j\ge0}\mathcal U_j(v)h^{-j}.
\]
Then, as a formal power series in $z$,
\begin{equation}\label{us:eq:forward-recipe}
 \sum_{j\ge0}\Theta_j(L,v;a)z^j
 =\exp\!\left(L\sum_{k\ge1}z^k d_k(\mathcal D;a)\right)
       \sum_{k\ge0}\mathcal U_k(v)z^k.
\end{equation}
All operations truncate after finitely many terms at a requested
order. For example
$\mathcal U_1=5v^2/6-2v$ and
$\mathcal U_2=25v^4/72-8v^3/3+4v^2$.
The exact gamma cumulant identity behind the compiler is
\[
 \log\mathbb E e^{t(Y_p-L)}
 =-Lt+hL\{\log(1+a/h)-\log(1+(a-t)/h)\}.
\]
Thus the forward coefficient rule is algebraic, and its uniform
error comes from bounded kernel derivatives and centered gamma moments.

\subsection{Formal inversion and a rigorous error bracket}

\begin{proof}[Proof of Theorem~\ref{us:thm:inverse}]
Set $z=1/h$ and write $p/h=L_0+\delta$, so $v=\lambda e^{-\delta}$.
The formal equation for the inverse is
\begin{equation}\label{us:eq:implicit}
 0=\lambda-\lambda e^{-\delta}
   +\log\!\left(
        \sum_{j\ge0}z^j
          \Theta_j(L_0+\delta,\lambda e^{-\delta};a)\right).
\end{equation}
It has a unique solution
$\delta=\sum_{j\ge0}z^{j+1}Q_j(L_0;\lambda,a)$ because its
coefficient of $\delta$ at $z=0,\delta=0$ is $\lambda\ne0$.
The recursion is effective: substitute known coefficients and divide
the next residual coefficient by $-\lambda$.

There are no hidden denominators in $\lambda$. For $j>0$ every
$\Theta_j(L,v;a)$ is divisible by $v$: this follows directly
from \eqref{us:eq:forward-recipe}, since $\mathcal D$ preserves
the ideal $(v)$, $d_k(0;a)=0$, and $\mathcal U_j(0)=0$ for $j>0$.
The logarithm in
\eqref{us:eq:implicit} is therefore coefficientwise divisible by
$\lambda$. After division the derivative with respect to $\delta$
is one. Induction gives $Q_j\in\mathbb Q[L_0,\lambda,a]$.
The bound on the degree in $L_0$ follows by the same induction:
a term of order $z^r$ has degree at most $r$ in $L_0$, including
products and the substitutions in \eqref{us:eq:implicit}.

At order $z$, the equation is
$\lambda Q_0+\Theta_1(L_0,\lambda;a)=0$, proving
\eqref{us:eq:Q0}. At order $z^2$ it gives
\[
 Q_1=\frac{Q_0^2}{2}
  -\frac{
     Q_0(\partial_L-\lambda\partial_\lambda)\Theta_1
       +\Theta_2-\Theta_1^2/2}{\lambda},
\]
evaluated at $(L_0,\lambda;a)$. Expanding the forward compiler
through order two gives \eqref{us:eq:Q1}. The accompanying exact
program checks these cancellations both in the logarithmic equation
and in the original normalized expression. A second program derives
the same answer directly from gamma cumulants after shifting $p$.

For the analytic error, let
$p_* = hL_0+\sum_{j=0}^Jh^{-j}Q_j$ and use
\eqref{us:eq:forward} through order $J+1$. The formal
cancellations prove
\[
 R_{p_*,h}(a)-\theta
 =O\!\left(h^{-J-2}L_0^{J+2}\right)
\]
uniformly on the stated compact sets. Apply the same finite expansion
at $p_*\pm K h^{-J-1}L_0^{J+2}$, with $K>0$ fixed.
The leading term changes by
\[
 \pm\theta\lambda K h^{-J-2}L_0^{J+2}(1+o(1)),
\]
and the change of the finitely many correction terms is of smaller
relative order. Since $\theta\lambda$ is bounded away from zero
on a compact subinterval of $(0,1)$, taking $K$ sufficiently large
places the two expanded values on opposite sides of $\theta$.
Both exponents remain in the same compact transition window.
Lemma~\ref{us:lem:monotone} then brackets the unique inverse
between them, proving \eqref{us:eq:inverse-all}. No derivative of
the remainder in \eqref{us:eq:forward} has been assumed.
\end{proof}

\subsection{A quantitative example}

For $\theta=1/2$ and $a=1/2$, numerical inversion of the exact
positive integral gives the values in Table~\ref{us:tab:inverse}.
The final three columns report absolute exponent errors, successively
using $hL_0$, $hL_0+Q_0$, and $hL_0+Q_0+Q_1/h$.
These are high-precision diagnostics of the defining integral,
not interval-certified root enclosures.

\begin{table}[htbp]
\centering
\caption{The inverse harmonic transition at $\theta=a=1/2$.}
\label{us:tab:inverse}
\begin{tabular}{@{}rrrrr@{}}
\toprule
$h$ & numerical $p_{1/2}$ & leading error & with $Q_0$ & with $Q_1/h$\\
\midrule
50 & $182.994297936847$ & $3.72486$ & $4.03044\cdot10^{-2}$ & $3.19479\cdot10^{-3}$\\
200 & $998.989394784027$ & $4.65277$ & $1.82329\cdot10^{-2}$ & $4.26953\cdot10^{-4}$\\
800 & $5091.951714446499$ & $5.56974$ & $7.10744\cdot10^{-3}$ & $4.78128\cdot10^{-5}$\\
\bottomrule
\end{tabular}
\end{table}

The increasing error of the leading approximation is expected:
$Q_0$ contains a multiple of $\log h$. The decreasing errors after
correction illustrate why solving only the leading transition profile
can be inadequate even when the forward normalized sum is well
approximated.
